\section{Практическая часть}

Перед измерением убедимся, что экран находится в зоне Фраунгофера:
\[
X = 283~\text{см} \ll \frac{\max \{l_1^2,\, l_2^2\}}{\lambda} = \frac{524.41~\text{см}^2}{3.3~\text{см}} \sim 159~\text{см}
\]

\subsection{Поле при отсутствии принимаемого сигнала}

\subsection{Второй метод измерения КНД}

При изменении положения антеннты относительно отражающего щита $X + \Delta X$ были измерены $|E_{max}|^2$ и $|E_{min}|^2$ в измерительной линии. Из измерений можно выразить КБВ в линии $\kappa = E_{min} / E_{max}$, откуда рассчитать коэффициент $\tilde{\Gamma}$:
\[
\tilde{\Gamma}(\Delta X) = \frac{1 - \kappa(\Delta X)}{1 + \kappa(\Delta X)} = \frac{1 - \sqrt{K(\Delta X)}}{1 + \sqrt{K(\Delta X)}}
\]

Полученная зависимость приведена на \cref{fig:ph:3}: